%=========================================================================
%  Deep Bayesian Active Learning with Image Data -- reproduction report
%  Compiles on Overleaf with pdfLaTeX out of the box (figures optional).
%  To add your plots: drop the PNGs into the figures/ folder using the
%  filenames referenced below (they match the repo's plotting scripts).
%=========================================================================
\documentclass[11pt]{article}

\usepackage[utf8]{inputenc}
\usepackage[T1]{fontenc}
\usepackage[margin=1in]{geometry}
\usepackage{amsmath,amssymb,amsfonts}
\usepackage{graphicx}
\usepackage{booktabs}
\usepackage{subcaption}
\usepackage{caption}
\usepackage{xcolor}
\usepackage{listings}
\usepackage[hidelinks]{hyperref}

\graphicspath{{figures/}}

% --- Placeholder helper: show the figure if present, else a grey box. ---
% (Filenames are only used as paths, never typeset, so underscores are safe.)
\newcommand{\figorbox}[3]{%
  \IfFileExists{figures/#1}%
    {\includegraphics[width=#2]{#1}}%
    {\fcolorbox{gray}{gray!10}{\parbox[c][3.2cm][c]{#2}{\centering\ttfamily
      \small [ figure placeholder ]\\[6pt] \footnotesize generate with:\\[2pt] #3}}}%
}

% --- Minimal Python listing style. ---
\definecolor{codegray}{gray}{0.95}
\lstset{
  basicstyle=\ttfamily\footnotesize,
  backgroundcolor=\color{codegray},
  keywordstyle=\color{blue!60!black},
  commentstyle=\color{green!40!black},
  stringstyle=\color{red!60!black},
  language=Python, breaklines=true, frame=single, framesep=4pt,
  numbers=left, numberstyle=\tiny\color{gray}, showstringspaces=false,
}

%=========================================================================
\title{\textbf{Deep Bayesian Active Learning with Image Data}\\[4pt]
       \large A Reproduction and Implementation Study}
\author{Your Name \\ \small Course: ASI \quad\textbar\quad
        \href{https://github.com/CaoNam04/ASI-deep-bayesian-active-learning}%
        {github.com/CaoNam04/ASI-deep-bayesian-active-learning}}
\date{\today}

\begin{document}
\maketitle

\begin{abstract}
% TODO: 4--6 sentence summary of what you did and what you found.
We reproduce the active-learning framework of Gal, Islam and Ghahramani
\cite{gal2017deepbal}, which combines Bayesian convolutional neural networks
with uncertainty-based acquisition functions to label image data efficiently.
We re-implement the method from scratch in PyTorch, evaluate five acquisition
functions on MNIST, study the role of model uncertainty by comparing a Bayesian
CNN to a deterministic one, compare against semi-supervised baselines, and apply
the method to melanoma diagnosis on the ISIC~2016 dataset. % TODO: headline result.
\end{abstract}

%-------------------------------------------------------------------------
\section{Introduction}
\label{sec:intro}
% TODO: motivate the problem in your own words.
Labelling data is often the costliest step in deploying a machine-learning
system, especially in domains such as medical imaging where labels require
expert time. \emph{Active learning} (AL) reduces this cost by letting the model
choose which unlabelled points are most worth labelling. Classical AL struggles
with high-dimensional image data and with models that do not express
uncertainty. The paper we reproduce addresses both issues by using a Bayesian
CNN whose predictive uncertainty drives the acquisition of new labels.

\paragraph{Contributions of this report.}
\begin{itemize}
  \item A clean, from-scratch PyTorch re-implementation of the method
        (Section~\ref{sec:method}--\ref{sec:impl}).
  \item Reproduction of the MNIST acquisition-function comparison
        (Section~\ref{sec:mnist}) and the Bayesian-vs-deterministic study
        (Section~\ref{sec:uncertainty}).
  \item A comparison to semi-supervised baselines (Section~\ref{sec:semisup})
        and an application to ISIC~2016 melanoma diagnosis
        (Section~\ref{sec:isic}).
  \item A discussion of implementation choices and differences from the original
        results (Section~\ref{sec:discussion}).
\end{itemize}

%-------------------------------------------------------------------------
\section{Background and Method}
\label{sec:method}

\subsection{The active-learning loop}
Given a small labelled set $\mathcal{D}_{\text{train}}$, a large unlabelled pool
$\mathcal{D}_{\text{pool}}$, and a model $\mathcal{M}$, an \emph{acquisition
function} $a(x,\mathcal{M})$ scores each pool point. At every step we select
\begin{equation}
  x^\star = \arg\max_{x \in \mathcal{D}_{\text{pool}}} a(x,\mathcal{M}),
\end{equation}
query its label from an oracle, add it to $\mathcal{D}_{\text{train}}$, and
retrain. The loop repeats, growing the labelled set over time.

\subsection{Bayesian CNN via MC dropout}
A Bayesian CNN places a distribution over weights $\omega$. Exact inference is
intractable, so we approximate the posterior $p(\omega\mid\mathcal{D})$ with a
dropout distribution $q^\star_\theta(\omega)$ and use \emph{MC dropout}: keeping
dropout active at test time and averaging $T$ stochastic forward passes
\cite{gal2016dropout},
\begin{equation}
  p(y=c\mid x,\mathcal{D}) \approx
  \frac{1}{T}\sum_{t=1}^{T} p(y=c\mid x,\hat{\omega}_t),
  \qquad \hat{\omega}_t \sim q^\star_\theta(\omega).
  \label{eq:mcdropout}
\end{equation}
A \emph{deterministic} CNN instead sets $q^\star_\theta(\omega)=\delta(\omega-\theta)$
(a point mass): it captures aleatoric (data) uncertainty but not epistemic
(model) uncertainty.

\subsection{Acquisition functions}
\label{sec:acq}
Let $\hat{p}^t_c$ be the probability of class $c$ on the $t$-th MC pass.

\paragraph{Max Entropy \cite{shannon1948}.}
\begin{equation}
  \mathbb{H}[y\mid x,\mathcal{D}] = -\sum_c p(y=c\mid x,\mathcal{D})\,
  \log p(y=c\mid x,\mathcal{D}).
\end{equation}

\paragraph{BALD \cite{houlsby2011bald}.}
Mutual information between prediction and weights, i.e.\ points on which the
sampled models \emph{disagree}:
\begin{align}
  \mathbb{I}[y,\omega\mid x,\mathcal{D}]
  &= \mathbb{H}[y\mid x,\mathcal{D}]
   - \mathbb{E}_{p(\omega\mid\mathcal{D})}\!\big[\mathbb{H}[y\mid x,\omega]\big]\\
  &\approx -\sum_c\Big(\tfrac{1}{T}\textstyle\sum_t \hat{p}^t_c\Big)
            \log\Big(\tfrac{1}{T}\textstyle\sum_t \hat{p}^t_c\Big)
          + \tfrac{1}{T}\sum_{c,t}\hat{p}^t_c\log\hat{p}^t_c .
  \label{eq:bald}
\end{align}

\paragraph{Variation Ratios \cite{freeman1965}.}
\begin{equation}
  \text{variation-ratio}[x] = 1 - \max_y p(y\mid x,\mathcal{D}).
\end{equation}

\paragraph{Mean STD.}
\begin{equation}
  \sigma_c = \sqrt{\mathbb{E}_q[p(y=c\mid x,\omega)^2]
                  - \mathbb{E}_q[p(y=c\mid x,\omega)]^2}, \qquad
  \sigma(x) = \frac{1}{C}\sum_c \sigma_c .
\end{equation}

\paragraph{Random.} Uniform selection from the pool (baseline).

%-------------------------------------------------------------------------
\section{Implementation}
\label{sec:impl}

\subsection{Model architecture}
Following the paper, the MNIST model is
\texttt{conv--relu--conv--relu--maxpool--dropout--dense--relu--dropout--dense},
with $32$ convolution kernels ($4\!\times\!4$), $2\!\times\!2$ pooling, a
$128$-unit dense layer, and dropout probabilities $0.25$ and $0.5$. For ISIC we
fine-tune VGG16 \cite{simonyan2015vgg} pre-trained on ImageNet, replacing the
$1000$-way head with a $2$-way (benign/malignant) head and keeping the two
$0.5$ dropout layers for MC dropout.

\subsection{Training and acquisition details}
% TODO: adjust if you changed any hyperparameters.
Each acquisition step trains a freshly initialised model to convergence and
keeps the weights with the best validation accuracy; the model is reset every
step to isolate the effect of the acquisition function. On MNIST we start from
$20$ balanced labelled points, acquire $10$ points per step using $T$ MC
samples, and evaluate on the standard $10$k test set.

\subsection{Code structure}
The implementation is split into modules; the core acquisition functions read as
follows (abridged):
\begin{lstlisting}
def bald(probs):  # probs: (N, T, C)
    mean_p = probs.mean(axis=1)
    entropy_of_mean = -(mean_p * np.log(mean_p + EPS)).sum(axis=1)
    entropy_per_pass = -(probs * np.log(probs + EPS)).sum(axis=2)
    return entropy_of_mean - entropy_per_pass.mean(axis=1)
\end{lstlisting}

%-------------------------------------------------------------------------
\section{Experiments and Results}
\label{sec:results}

\subsection{MNIST: comparison of acquisition functions}
\label{sec:mnist}
% TODO: describe what you observe in your run.
\begin{figure}[t]
  \centering
  \figorbox{accuracy_comparison.png}{0.8\linewidth}%
           {python plot\_results.py}
  \caption{MNIST test accuracy vs.\ number of acquired images, averaged over
           repetitions (shaded: std). \emph{TODO: comment on which functions
           dominate.}}
  \label{fig:mnist}
\end{figure}

\begin{table}[t]
  \centering
  \caption{Number of acquired images needed to reach a given test error.
           \emph{TODO: fill in from your runs.}}
  \label{tab:acq-steps}
  \begin{tabular}{lccccc}
    \toprule
    \% error & BALD & Var Ratios & Max Ent & Mean STD & Random\\
    \midrule
    10\% & --- & --- & --- & --- & ---\\
    5\%  & --- & --- & --- & --- & ---\\
    \bottomrule
  \end{tabular}
\end{table}

\subsection{Importance of model uncertainty}
\label{sec:uncertainty}
% TODO: comment on the gap between red and blue.
\begin{figure}[t]
  \centering
  \figorbox{figure2_bayesian_vs_deterministic.png}{\linewidth}%
           {python plot\_figure2.py}
  \caption{Bayesian CNN (red) vs.\ deterministic CNN (blue) for BALD, Variation
           Ratios and Max Entropy. The deterministic model lacks epistemic
           uncertainty, so BALD degenerates toward random selection.}
  \label{fig:figure2}
\end{figure}

\subsection{Comparison to semi-supervised learning}
\label{sec:semisup}
\begin{table}[t]
  \centering
  \caption{Test error on MNIST with $1000$ labelled images. Semi-supervised
           methods additionally use $\sim$49k--59k unlabelled images; our active
           learner uses only the $1000$ acquired labels. \emph{TODO: fill in your
           active-learning errors via} \texttt{report\_table.py}.}
  \label{tab:semisup}
  \begin{tabular}{lc}
    \toprule
    Method & Test error\\
    \midrule
    \multicolumn{2}{l}{\emph{Semi-supervised (from \cite{gal2017deepbal})}}\\
    DGN \cite{kingma2014dgn}                 & 2.40\%\\
    Ladder Network ($\Gamma$-model)          & 1.53\%\\
    Ladder Network (full)                    & 0.84\%\\
    \midrule
    \multicolumn{2}{l}{\emph{Active learning (this work)}}\\
    Random       & ---\\
    BALD         & ---\\
    Max Entropy  & ---\\
    Var Ratios   & ---\\
    \bottomrule
  \end{tabular}
\end{table}

\subsection{ISIC 2016 melanoma diagnosis}
\label{sec:isic}
% TODO: describe AUC behaviour and positive-acquisition behaviour.
\begin{figure}[t]
  \centering
  \figorbox{isic_figure5.png}{\linewidth}%
           {python -m isic.plot\_results}
  \caption{ISIC~2016: test AUC (left) and number of acquired positive
           (malignant) examples (right) for BALD vs.\ uniform acquisition, on
           two random test splits. \emph{TODO: comment on BALD acquiring more
           positives and reaching higher AUC.}}
  \label{fig:isic}
\end{figure}

%-------------------------------------------------------------------------
\section{Discussion}
\label{sec:discussion}
% TODO: discuss differences from the paper and design choices.
\begin{itemize}
  \item \textbf{Training stability.} Reporting the best-validation checkpoint
        (rather than the final epoch) removes sharp accuracy dips caused by
        retraining from scratch each step.
  \item \textbf{Variation Ratios.} We use the soft definition
        $1-\max_c p(y=c\mid x)$, which matches the paper's formula and is also
        well defined for a single deterministic pass.
  \item \textbf{Differences from the paper.} % TODO: e.g. accuracy slightly
        higher in the mid-curve; Mean STD variance with few repetitions, etc.
\end{itemize}

%-------------------------------------------------------------------------
\section{Conclusion}
\label{sec:conclusion}
% TODO: 3--4 sentences.
We re-implemented and reproduced the main results of \cite{gal2017deepbal}.
Uncertainty-aware acquisition (BALD, Variation Ratios, Max Entropy) substantially
reduces the number of labels needed compared to random sampling, and the
epistemic uncertainty captured by a Bayesian CNN is essential to this gain.

%-------------------------------------------------------------------------
\begin{thebibliography}{9}
\bibitem{gal2017deepbal} Y.~Gal, R.~Islam, Z.~Ghahramani.
  \emph{Deep Bayesian Active Learning with Image Data}. ICML, 2017.
\bibitem{gal2016dropout} Y.~Gal, Z.~Ghahramani.
  \emph{Dropout as a Bayesian Approximation}. ICML, 2016.
\bibitem{houlsby2011bald} N.~Houlsby, F.~Husz\'ar, Z.~Ghahramani, M.~Lengyel.
  \emph{Bayesian Active Learning for Classification and Preference Learning}.
  arXiv:1112.5745, 2011.
\bibitem{shannon1948} C.~E.~Shannon.
  \emph{A Mathematical Theory of Communication}. Bell System Tech.\ J., 1948.
\bibitem{freeman1965} L.~G.~Freeman. \emph{Elementary Applied Statistics}. 1965.
\bibitem{simonyan2015vgg} K.~Simonyan, A.~Zisserman.
  \emph{Very Deep Convolutional Networks for Large-Scale Image Recognition}.
  ICLR, 2015.
\bibitem{kingma2014dgn} D.~P.~Kingma et al.
  \emph{Semi-supervised Learning with Deep Generative Models}. NIPS, 2014.
\bibitem{gutman2016isic} D.~Gutman et al.
  \emph{Skin Lesion Analysis toward Melanoma Detection (ISBI 2016 / ISIC)}.
  arXiv:1605.01397, 2016.
\end{thebibliography}

\end{document}
